% !TeX root = ../数值分析与科学计算笔记.tex
% 由主文档通过 \input 载入；请编译主文档。
\setcounter{chapter}{1} % 保留原书章号。
\chapter{公式使用前的条件核对}
\label{app:na-conditions}

\section{存在性、误差与计算保证的区别}
求根中的连续异号保证根存在，导数非零保证局部 Newton 行为，两者承担不同作用；最小二乘解存在不意味着唯一，列满秩才保证参数唯一；正定性保证 Cholesky 和 CG 的核心结构，却不意味着任意步长的时间积分稳定。检查条件时，应将所使用定理的对象与实际输入逐项对应。

误差界中的常数常依赖导数上界、区间长度或条件数。一个形式上为 $O(h^4)$ 的方法，在四阶导数很大时可能不如常数较小的二阶方法；如果函数不满足四阶光滑条件，则该阶数甚至没有适用依据。局部缺陷、全局误差、代数残差和随机标准误差分别描述不同量，不应在图表或报告中只写一个未定义的“误差”。

\section{正文未完整展开的理论}
为保持本科主线，正文对奇异值分解给出了基于谱定理的证明要点，对 CG 给出投影与多项式误差界的逻辑，对割线法阶数、一般 Runge--Kutta 高阶条件、线性多步法等价定理和一般 QR 迭代仅说明结论或证明范围。进一步学习应专门补足这些证明及浮点误差分析。正文中已展开的基础证明包括压缩迭代、Newton 局部误差、线性系统扰动、Cholesky、最小二乘正规条件、插值余项、Gauss 精确性、Euler 全局界和一维差分稳定收敛。

